\subsection{ML\_dEdge: Sequence-to-Property Prediction Model}

The ML\_dEdge model employs a Transformer-based encoder architecture to predict the difference in edge statistics (dEdge) between dimer and monomer peptide simulations from amino acid sequences.
The model architecture consists of a 6-layer Transformer encoder with 8-head self-attention, 512-dimensional embeddings, and 2048-dimensional feed-forward networks.
The encoder processes peptide sequences of variable length (up to 10 amino acids) and outputs a continuous dEdge value through a regression head.
Training is performed using curriculum learning across two iterations: the first iteration establishes baseline performance on a comprehensive dataset of cysteine-containing peptides (dipeptides, tripeptides, and tetrapeptides), while the second iteration introduces additional sequence variations through a three-stage curriculum learning approach with progressively decreasing learning rates (0.2, 0.1, 0.05) and warmup epochs.
The model is trained to minimize the mean squared error (MSE) loss between predicted and actual dEdge values:
\begin{equation}
L_{\text{MSE}} = \frac{1}{N} \sum_{i=1}^{N} (y_i - \hat{y}_i)^2
\end{equation}
where $N$ is the batch size, $y_i$ is the true dEdge value for sequence $i$, and $\hat{y}_i$ is the predicted dEdge value from the model.
Training is performed using the Adam optimizer with gradient clipping for stability.
The final model achieves strong predictive performance, enabling accurate estimation of peptide self-assembly properties from sequence information alone.

\subsection{ML\_dEdge\_gen: Conditional Generative Model for Sequence Design}

The ML\_dEdge\_gen model implements a conditional Transformer decoder architecture for generating peptide sequences that match target dEdge values and sequence lengths.
Unlike the ML\_dEdge model which predicts properties from sequences, this generative model takes target dEdge value ranges and sequence length specifications as input and generates corresponding peptide sequences.
The architecture consists of a decoder-only Transformer with 8 layers, 12 attention heads, 768-dimensional embeddings, and 3072-dimensional feed-forward networks.
Conditional inputs (dEdge value and sequence length) are embedded and used as encoder-like context throughout the decoder layers, enabling the model to generate sequences autoregressively while being conditioned on target properties.
Training employs adversarial learning with a frozen pre-trained ML\_dEdge model as a critic, combined with soft reconstruction loss for sequence learning.
Training uses scheduled sampling to gradually transition from teacher forcing (100\% ground truth tokens) to mixed sampling (50\% teacher forcing, 50\% generated tokens) over epochs, and temperature annealing that starts at 2.0 for exploration and decreases to 1.0 for focused generation.
The training loss consists of two components: (1) a soft reconstruction loss $L_{\text{recon}}$ that learns sequence structure from similar sequences with similar dEdge values:
\begin{equation}
L_{\text{recon}} = \frac{1}{N} \sum_{i=1}^{N} \text{KL}(P_{\text{predicted},i} || P_{\text{soft},i})
\end{equation}
where $N$ is the batch size, $P_{\text{predicted},i}$ is the predicted token distribution for sequence $i$, and $P_{\text{soft},i}$ is a soft target distribution created from a weighted average of similar sequences (sequences with dEdge values within a threshold of the target dEdge).
The soft target is constructed by finding the top-$k$ most similar sequences (by dEdge value) and creating a weighted probability distribution over the vocabulary at each position, where weights are inversely proportional to dEdge difference.
This approach enables the model to learn from multiple valid sequences that map to similar dEdge values, rather than forcing exact sequence matching.
(2) A dEdge matching loss $L_{\text{dedge}}$ that ensures generated sequences match target dEdge values:
\begin{equation}
L_{\text{dedge}} = \frac{1}{N} \sum_{i=1}^{N} (d_{\text{target},i} - d_{\text{predicted},i})^2
\end{equation}
where $d_{\text{target},i}$ is the target dEdge value and $d_{\text{predicted},i}$ is the dEdge value predicted by the frozen ML\_dEdge critic model for the generated sequence.
The total training loss combines these components, with the reconstruction loss scaled by vocabulary size to balance the loss magnitudes:
\begin{equation}
L_{\text{total}} = L_{\text{dedge}} + \frac{L_{\text{recon}}}{|\mathcal{V}|}
\end{equation}
where $|\mathcal{V}|$ is the vocabulary size (21 tokens).
Sequence generation uses temperature sampling with linear annealing from $T = 2.0$ (initial) to $T = 1.0$ (final) to balance exploration and exploitation during training.
Scheduled sampling is employed to gradually reduce teacher forcing probability from 1.0 (initial) to 0.5 (final) over training epochs, helping bridge the gap between training and inference.
During validation, the dEdge matching loss is computed to assess property matching, with this loss used for model selection.
To ensure proper evaluation of generalization, the training data uses stratified splits by $(dEdge, \text{length})$ combinations, ensuring that validation and test sets contain condition combinations not seen during training.
This approach enables the model to generate novel peptide sequences that match specified self-assembly properties, facilitating inverse design of peptides with desired characteristics.

\subsection{NHM\_DL: Two-Stage Network Hamiltonian Model Parameter Prediction}

The NHM\_DL model employs a two-stage Transformer-based architecture to predict Network Hamiltonian Model (NHM) parameters from peptide sequences.
Unlike standard regression approaches, this model addresses the inherent sparsity of NHM parameters, where many parameters are zero (not selected) for most sequences.
The architecture consists of a shared Transformer encoder (6 layers, 8 attention heads, 512-dimensional embeddings, 2048-dimensional feed-forward networks) that processes peptide sequences, followed by two parallel prediction heads.

The first stage performs binary classification to predict which NHM parameters are selected (non-zero) for a given sequence.
This selection head outputs binary logits for each parameter, indicating the probability that the parameter is relevant for the input sequence.
The selection loss $L_{\text{selection}}$ uses binary cross-entropy:
\begin{equation}
L_{\text{selection}} = \frac{1}{N} \sum_{i=1}^{N} \sum_{j=1}^{P} \text{BCE}(s_{ij}, \hat{s}_{ij})
\end{equation}
where $N$ is the batch size, $P$ is the number of NHM parameters, $s_{ij}$ is the true binary selection mask (1 if parameter $j$ is non-zero for sequence $i$, 0 otherwise), and $\hat{s}_{ij}$ is the predicted selection logit.

The second stage performs regression to predict the actual values of selected parameters.
The regression head outputs continuous values for all parameters, but the loss is masked to only penalize errors for parameters that are actually selected.
The masked regression loss $L_{\text{regression}}$ is computed as:
\begin{equation}
L_{\text{regression}} = \frac{1}{N} \sum_{i=1}^{N} \frac{\sum_{j=1}^{P} s_{ij} \cdot (y_{ij} - \hat{y}_{ij})^2}{\sum_{j=1}^{P} s_{ij}}
\end{equation}
where $y_{ij}$ is the true value of parameter $j$ for sequence $i$, $\hat{y}_{ij}$ is the predicted value, and the denominator normalizes by the number of selected parameters to avoid bias from sequences with different numbers of active parameters.

The total training loss combines both stages:
\begin{equation}
L_{\text{total}} = L_{\text{selection}} + L_{\text{regression}}
\end{equation}
This two-stage approach allows the model to first learn which parameters are relevant for each sequence type, then focus on accurately predicting the values of those selected parameters.
Training uses the Adam optimizer with learning rate scheduling, and data is split into training (70\%), validation (10\%), and test (20\%) sets.
NHM parameters are normalized using per-parameter MinMaxScaler to ensure consistent scaling across different parameter types.
Separate models are trained for monomer and dimer states, as the NHM parameter distributions differ significantly between these states.
The model enables efficient prediction of complex network statistics from sequence information alone, facilitating high-throughput analysis of peptide self-assembly properties.

